\section{Scope and Systematization Method}\label{sec:methodology}\label{sec:background}
\subsection{Research questions, eligibility and threat polarity}
We organize the systematization around four questions. \textbf{RQ1 (landscape):} which hidden targets, threat polarities and auxiliary observables appear across LLM and agent systems, and where is each signal generated versus observed? \textbf{RQ2 (comparability):} when do two reported defense results measure the same target, receiver, intervention and cost object? \textbf{RQ3 (joint evidence):} which results bind protection, utility and cost at one evaluated operating point? \textbf{RQ4 (transfer):} which evidence distinguishes attack generalization, defended-condition evaluation and defense transfer across tasks, models, topologies and hardware? RQ1 supplies the field context; RQ2--RQ4 correspond to the three contributions.

Core inclusion requires five conjunctive criteria: (C1) LLM or agent execution, (C2) a hidden target property, (C3) an execution-dependent auxiliary observable rather than ordinary output semantics alone, (C4) specified observer permissions, and (C5) an empirical inference evaluation. Mixed-feature papers qualify only if they isolate an auxiliary-only C3 component; thus coding fingerprints, WildAuth, ZeroPrint, endpoint-content fingerprints and StepFinder remain adjacent boundary comparators \cite{ghaleb2026codingfingerprint,wang2026dialogue,shao2026zeroprint,leshin2026endpoint,zhu2026stepfinder}. Table~\ref{tab:boundary-rules} states the observable eligibility rules. For C3, the extraction must name an execution mechanism and an evaluated observable separate from ordinary linguistic meaning. A filter qualifies when its accept/reject behavior or differential response is the measured oracle for a hidden property; merely analyzing a reply's wording or style does not. If a mixed paper does not isolate that mechanism, its membership stays adjacent. This operational distinction concerns the measured signal, not whether the output is text.
\begin{table}[t]\centering\footnotesize
\caption{Observable-based eligibility. A hidden target alone does not make a semantic attack a side channel.}\label{tab:boundary-rules}
\begin{tabularx}{\columnwidth}{@{}p{.39\columnwidth}Y@{}}\toprule
Observed input & Decision and required distinction\\\midrule
Length, timing, cache response & Core when measured execution dependence supports hidden-target inference.\\
Own-site actions & Core attribution with site permission; no automatic access to other sites or private tool logs.\\
Reply states the secret & Direct disclosure; adjacent, not auxiliary-channel evidence.\\
Reply style/content identifies a user & Adjacent unless a separately evaluated auxiliary mechanism is established. WildAuth stays a mixed-view comparator.\\
Output-filter membership oracle & Can be core: the filter's execution, not ordinary semantic output inference, mediates the signal.\\
Privileged loop/fault logs & Separate diagnostic component when C1--C5 hold; does not increase confidentiality evidence.\\\bottomrule
\end{tabularx}\end{table}


We use four narrative groupings of objectives and observer roles, not mutually exclusive classes. An unauthorized identity inference can be both attribution and a confidentiality concern. Each record separately specifies the inference target, observer, authorization relation and protected interest:
\textbf{Type I} concerns confidentiality/privacy; \textbf{Type II} concerns model, harness or identity attribution; \textbf{Type III} concerns evasion of a platform's human--agent detector; \textbf{Type IV} concerns authorized operator diagnostics. Types can overlap. WRA protects research intent (I); UI Traces attributes agents (II); AHB evades a defensive detector (III); ClawGuard diagnoses loops (IV) \cite{jeong2026wra,uitraces2026,zhu2026turing,gan2026clawguard}. These roles are mapped per core family, alongside generating and receiving layers, in Table~\ref{tab:core-family-map}.

\subsection{Corpus construction and reconciliation}
Forty seeds, venue-targeted proceedings screening and one-hop citation snowballing produced an initial map of \SnapshotN{} families (\SnapshotCore{} core, \SnapshotAdjacent{} adjacent). A supplementary database and bulk keyword search returned 615 occurrences and 565 unique records; auditing its 204 exclusions rerouted four records, and the search added 15 core and seven adjacent families. The final collection has \FinalCore{} core and \FinalAdjacent{} adjacent families. The completed search frames, closed on October~21, 2026, combine Semantic Scholar and native arXiv queries with venue-targeted screening of S\&P, USENIX Security, NDSS, CCS and PMLR. C1--C5 screening and version deduplication operate on records before family reconciliation; the supplement retains queries, dispositions and excluded incomplete search frames. Architecture and hardware venues were not comprehensively screened; the finite frame is not an exhaustive census.

\subsection{Source-bound result records}
Each record links its target, observation boundary and operating point to a source version and section/table/figure locator. The authors report reviewing the classification of all 82 families using titles, abstracts and experimental sections, and checking all 39 result records and 180 family-audit cells against source locators. The subsequent 351-cell record audit and 133 within-question pair decisions retain author-assigned reasons; executable checks establish their accounting consistency. Candidate fields produced with AI assistance were subject to this author review.

Appendices~\ref{app:records} and \ref{app:reconstruction} contain the record index, comparison summary and query memberships in this PDF. The ledger also separates cited publication status, inspected version, measurement mode and remaining binding uncertainty. Published status increases traceability of venue review but does not establish the correctness of a result; a preprint's observations remain provisional at their inspected scope.

\soktakeaway{SoK Takeaway 1 (Map evidence to the receiving deployment)}{A family-level label does not identify which signal, observer permission or protected objective a result supports. The complete core map separates signal origin from receiving access; the result ledger then binds each comparison to a source version and operating point.}
